\section{Results}

\subsection{Within-Dataset Performance}
Table~\ref{tab:within} summarizes within-dataset performance as fold-mean AUC $\pm$ SD under the leakage-clean protocol (LOCO on CALCE, 5-fold grouped CV on NASA), with Platt Brier scores. Every tree configuration clears 0.89 at every horizon on both datasets (Fig.~\ref{fig:within}); the GRU matches the means with higher fold variance, and the single-fit hazard model lands in the same range. On the four new Battery Archive groups, @@WITHIN_BA_TOTAL@@ of @@WITHIN_BA_COMBOS@@ tree--group--horizon configurations reach 0.85, spanning @@WITHIN_BA_RANGE@@ --- though the mostly-censored LFP-SNL group (@@BA_PREV_LFP@@ prevalence) is supporting evidence only.

\begin{table*}[!t]
\caption{Within-Dataset Failure-Risk Performance Across Four Horizons (Fold-Mean AUC $\pm$ Fold SD; Tree Ensembles and GRU Platt-Calibrated, Hazard Models Raw; CALCE Uses Leave-One-Cell-Out, NASA 5-Fold Cell-Grouped CV). Brier Is the Platt Pooled Value Averaged Over Horizons (Raw for Hazard Rows).}
\label{tab:within}
\centering
\resizebox{\textwidth}{!}{%
\begin{tabular}{llcccccc}
\toprule
\textbf{Model} & \textbf{Dataset} & \textbf{H=10} & \textbf{H=20} & \textbf{H=30} & \textbf{H=50} & \textbf{Mean AUC} & \textbf{Brier}\\
\midrule
@@TAB_WITHIN@@
\bottomrule
\end{tabular}}
\end{table*}

\begin{figure*}[!t]
\centering
\includegraphics[width=0.62\textwidth]{Fig01_Within_Dataset_AUC.png}
\caption{Within-dataset AUC heatmap for the tree ensembles, mean across horizons, Platt-calibrated. Right block: the four new Battery Archive groups. GRU and hazard within-dataset numbers are in Table~\ref{tab:within}.}
\label{fig:within}
\end{figure*}

\subsection{Calibration Under Cross-Fitting}
Table~\ref{tab:cal} compares the three calibrators, averaged over the tree models and horizons, where every calibrator is fitted on cross-fitted out-of-fold scores (Section~II-H). Three observations follow. First, on discrimination the raw scores are effectively unbeatable within dataset: any fitted calibrator is a monotone re-mapping and preserves per-fold ranking up to the ties it introduces, and the fold-mean AUC differences between calibrators are correspondingly small. Second, the step function of isotonic regression creates large tied groups whose cost is invisible to fold-mean ROC-AUC but severe for precision--recall: on NASA, raw scores reach a mean PR-AUC of @@NASA_PRAUC_RAW@@ while isotonic collapses to @@NASA_PRAUC_ISO@@ (Fig.~\ref{fig:prauc}). Third, calibration quality proper is best for Platt on CALCE (ECE @@CALCE_PLATT_ECE@@ against @@CALCE_ISO_ECE@@ for isotonic), but the fitted Platt slope on CALCE is only @@CALCE_PLATT_SLOPE@@: on that long-tailed, leave-one-cell-out setting the out-of-fold sigmoid essentially washes confidence out toward the base rate while preserving ranking. Temperature scaling is the most balanced compromise on NASA (slope @@NASA_TEMP_SLOPE@@, Brier @@NASA_TEMP_BRIER@@, both better than Platt's). There is no global winner across datasets: temperature leads on NASA, Platt on CALCE ECE at the cost of a destroyed slope. We therefore read the comparison as: \emph{Platt preserves discrimination best while achieving calibration error comparable to or better than isotonic}, which is a materially weaker and more defensible claim than ``Platt is the better calibrator,'' and raw scores remain the safest ranking signal whenever the operating context is uncertain.

\begin{table*}[!t]
\caption{Calibration Quality Averaged Over Tree Models and Horizons. All Calibrators Are Fitted on Cross-Fitted Out-of-Fold Scores. Slope 1 / Intercept 0 Indicate Perfect Probabilistic Calibration.}
\label{tab:cal}
\centering
\resizebox{\textwidth}{!}{%
\begin{tabular}{llcccccc}
\toprule
\textbf{Dataset} & \textbf{Method} & \textbf{AUC (fold mean)} & \textbf{Brier} & \textbf{ECE} & \textbf{Log loss} & \textbf{Slope} & \textbf{Intercept}\\
\midrule
@@TAB_CAL@@
\bottomrule
\end{tabular}}
\end{table*}

\begin{figure*}[!t]
\centering
\includegraphics[width=0.9\textwidth]{fig_reliability_v2.png}
\caption{Reliability diagrams under the cross-fitted protocol. (a) Within CALCE, XGBoost at $H{=}20$: the fitted Platt sigmoid preserves ranking but washes confidence toward the base rate. (b) ALL-LCO $\to$ Severson without SOH at $H{=}20$: no fitted calibrator, whichever source scores it is applied to, repairs calibration on the shifted target.}
\label{fig:cal}
\end{figure*}

\begin{figure}[!t]
\centering
\includegraphics[width=\columnwidth]{fig_prauc_horizon.png}
\caption{Within-dataset PR-AUC of raw scores versus horizon. Precision--recall performance, unlike ROC-AUC, is sensitive to the horizon-dependent failure prevalence; it is also where isotonic's tied step outputs pay for their ranking-free compression.}
\label{fig:prauc}
\end{figure}

\subsection{Cross-Dataset Transfer With and Without SOH}
Table~\ref{tab:crosswith} reports LCO$\to$LFP transfer at $H{=}20$ with SOH kept, and Table~\ref{tab:crossno} without; Fig.~\ref{fig:collapse} maps the AUC drop from removing SOH. Every number carries a cell-level bootstrap interval plus per-cell mean and standard deviation. With SOH, the trees reach pooled raw AUC up to @@XFER_WITH_MAX@@ (Oxford) and @@XFER_WITH_SEV@@ (Severson). Without SOH, transfer collapses to at most @@XFER_NO_MAX@@ on Oxford and @@XFER_NO_SEV_MIN@@--@@XFER_NO_SEV_MAX@@ on Severson --- point estimates fall below 0.5 in several configurations, significantly so only for the hazard model on Oxford.

The GRU shows a source-dependent instability at three seeds: per-cell AUC on Severson without SOH is $0.987\pm0.015$ for NASA (stable, near-oracle) but $0.479\pm0.451$ for CALCE (seeds 0.21--1.00) and $0.893\pm0.131$ pooled --- against $0.998\pm0.003$, $0.964\pm0.033$ and $0.991\pm0.015$ with SOH. On the identical common feature set the trees retain their table~\ref{tab:crosswith} behavior, so the gap is architectural, not featural. On Oxford the GRU is at or below chance (Table~\ref{tab:gru}).

\begin{table*}[!t]
\caption{Cross-Dataset Transfer at H=20 With SOH (Upper-Bound / Label-Proximity Regime). Raw-Score Pooled AUC With Cell-Level Bootstrap 95\% CI; Per-Cell Mean $\pm$ SD Across Test Cells (Oxford: Five Cells, One With a Single Positive Row; See Section II-B).}
\label{tab:crosswith}
\centering
\resizebox{\textwidth}{!}{%
\begin{tabular}{llcccc}
\toprule
 & & \multicolumn{2}{c}{\textbf{Oxford}} & \multicolumn{2}{c}{\textbf{Severson}}\\
\cmidrule(lr){3-4}\cmidrule(l){5-6}
\textbf{Training} & \textbf{Model} & \textbf{Pooled AUC [CI]} & \textbf{Per-cell} & \textbf{Pooled AUC [CI]} & \textbf{Per-cell}\\
\midrule
@@TAB_CROSS_WITH@@
\bottomrule
\end{tabular}}
\end{table*}

\begin{table*}[!t]
\caption{Cross-Dataset Transfer at H=20 Without SOH (Prognostically Meaningful Regime). Columns as in Table~\ref{tab:crosswith}.}
\label{tab:crossno}
\centering
\resizebox{\textwidth}{!}{%
\begin{tabular}{llcccc}
\toprule
 & & \multicolumn{2}{c}{\textbf{Oxford}} & \multicolumn{2}{c}{\textbf{Severson}}\\
\cmidrule(lr){3-4}\cmidrule(l){5-6}
\textbf{Training} & \textbf{Model} & \textbf{Pooled AUC [CI]} & \textbf{Per-cell} & \textbf{Pooled AUC [CI]} & \textbf{Per-cell}\\
\midrule
@@TAB_CROSS_NO@@
\bottomrule
\end{tabular}}
\end{table*}

\begin{figure*}[!t]
\centering
\includegraphics[width=0.95\textwidth]{fig_collapse_map.png}
\caption{Transferability collapse map at $H{=}20$ for the tree ensembles: the AUC drop from removing SOH, per source$\to$target pair and model. Positive values mean SOH removal destroys apparent transfer performance; the near-uniform saturation at the maximum drop is the signature of a single-feature shortcut rather than of distributed degradation knowledge. Right block: ALL-LCO $\to$ new Battery Archive targets. GRU deltas are reported in the text.}
\label{fig:collapse}
\end{figure*}

\begin{table}[!t]
\caption{GRU Transfer at H=20 on the Common Feature Set: Seed-Averaged Per-Cell AUC ($\pm$ SD Over Seeds 42, 1, 7).}
\label{tab:gru}
\centering
\small
\resizebox{\columnwidth}{!}{%
\begin{tabular}{llcc}
\toprule
\textbf{Training} & \textbf{Feature set} & \textbf{Oxford} & \textbf{Severson}\\
\midrule
@@TAB_GRU@@
\bottomrule
\end{tabular}}
\end{table}

The reading that survives the uncertainty treatment is deliberately modest. SOH appears on both sides of the learning problem --- in the feature vector and inside the label definition --- so high with-SOH transfer is compatible with two mechanisms: a chemistry-dependent mapping from SOH to remaining life, and mere detection of proximity to the SOH $=0.80$ criterion that defines the label. The experiments that follow separate these mechanisms.

Diagnosis or prognosis? Because the label window includes the scoring cycle, rows from already-failed cells score positive. Table~\ref{tab:prog} splits headline transfer numbers into all-rows AUC and prognosis-only AUC on the same scores restricted to rows with SOH still above 0.80. On Severson at $H{=}20$, XGBoost with SOH moves from @@PROG_SEV_ALL@@ to @@PROG_SEV@@ and without SOH from @@PROG_SEV_NO_ALL@@ to @@PROG_SEV_NO@@: 99\% of rows are scored above threshold, so concurrent detection explains only about 0.05--0.07 there; the rest is near-threshold ranking, priced by the failure-definition ablation in Section~III-F. The distance rule behaves identically (@@PROG_DIST_ALL@@ to @@PROG_DIST@@). On the new targets the diagnostic share is larger wherever cells linger below threshold --- HNEI @@PROG_HNEI@@, NCA @@PROG_NCA@@, NMC-SNL @@PROG_NMC@@, LFP-SNL @@PROG_LFP@@. The split is conservative --- rows just above 0.80 still count as prognosis --- so no with-SOH number here is pure prognosis: all are diagnostic-proximity upper bounds.

\begin{table*}[!t]
\caption{Diagnosis Versus Prognosis at H=20 (ALL-LCO Source, Raw Scores): All-Rows Pooled AUC [Cell-Level CI Where Computed] Against Prognosis-Only AUC on Rows Scored While SOH$>$0.80, With the Share of Rows Scored Above Threshold. GRU and Distance-Rule Rows Are Point Estimates.}
\label{tab:prog}
\centering
\resizebox{\textwidth}{!}{%
\begin{tabular}{lllccc}
\toprule
\textbf{Target} & \textbf{Model} & \textbf{Features} & \textbf{All rows} & \textbf{Prognosis-only} & \textbf{Frac.\ above}\\
\midrule
@@TAB_PROGNOSIS@@
\bottomrule
\end{tabular}}
\end{table*}

\subsection{Minimal Baselines: How Much Is the SOH Shortcut?}
Table~\ref{tab:baselines} places the tuned ensembles next to trivial baselines. The distance-to-threshold rule $d_t = 0.80 - \mathrm{SOH}_t$, which involves no fitting at all, transfers to Severson at pooled AUC @@BASE_DIST_ALL_SEV@@ [CI] for the ALL-LCO source --- matching or exceeding the tuned ensembles' with-SOH transfer in Table~\ref{tab:crosswith}. An SOH-only logistic regression behaves equivalently (@@BASE_SOH_ALL_SEV@@). Both are diagnostic baselines by construction: they order cells by proximity to the label threshold, not by learned degradation dynamics. By contrast, the sensor-only logistic regression transfers at @@BASE_SENSORS_NASA_SEV@@ on Severson (NASA source): below chance, meaning the remaining sensor features actively invert on the shifted distribution. Within CALCE, the cycle index alone ranks failure at pooled AUC @@BASE_CYCLE_CALCE@@, approaching the tuned ensembles, because CALCE's seven cells span very different cycle lives; on NASA, whose cells age under randomized profiles, the same feature is inverted (@@BASE_CYCLE_NASA@@).

Three conclusions follow. First, most of what transfers is the SOH coordinate, not degradation knowledge. Second, ``more features'' does not help transfer; sensor features carry dataset-specific artifacts that can invert. Third, the ensembles' within-dataset value is undiminished --- but the transfer claim must be downgraded from ``models transfer'' to ``the health state transfers, and models can read it.''

\begin{table*}[!t]
\caption{Minimal Transfer Baselines at H=20: Pooled AUC With Cell-Level Bootstrap 95\% CI. The Distance Rule Is the Unfitted Score $0.80-$SOH.}
\label{tab:baselines}
\centering
\resizebox{\textwidth}{!}{%
\begin{tabular}{lcccccc}
\toprule
 & \multicolumn{3}{c}{\textbf{Oxford}} & \multicolumn{3}{c}{\textbf{Severson}}\\
\cmidrule(lr){2-4}\cmidrule(l){5-7}
\textbf{Baseline} & \textbf{NASA src} & \textbf{CALCE src} & \textbf{ALL src} & \textbf{NASA src} & \textbf{CALCE src} & \textbf{ALL src}\\
\midrule
@@TAB_BASELINES@@
\bottomrule
\end{tabular}}
\end{table*}

\subsection{Factorial Feature Ablation}
Table~\ref{tab:abl} runs the full $2^3-1$ grid over \{SOH, cycle index, sensors\} at $H{=}20$. Three regularities stand out. SOH-only is the strongest single group on both LFP targets, confirming that the health coordinate dominates whatever transfer exists. Cycle-only is competitive on Oxford --- whose five cells fail at well-separated raw cycle counts, so the cycle index is a between-cell aging proxy --- but the pattern is dataset-specific and does not survive to Severson. No combination without SOH reaches 0.76 on either target (best @@ABL_MAX_NO_SOH@@), and removing both SOH and the cycle index leaves the sensors with essentially nothing (@@ABL_SENSORS_ONLY@@ with the ALL-LCO source on Severson). The ablation makes the shortcut diagnosis feature-by-feature: it is not that the models fail to use the sensor features; it is that the sensor features themselves carry no cross-dataset signal.

\begin{table*}[!t]
\caption{Factorial Feature Ablation at H=20 (XGBoost, Pooled AUC With Cell-Level Bootstrap 95\% CI). Feature Groups: SOH; Cycle Index; Sensors = Voltages, Current, Temperature, Duration.}
\label{tab:abl}
\centering
\resizebox{\textwidth}{!}{%
\begin{tabular}{lcccccc}
\toprule
 & \multicolumn{3}{c}{\textbf{Oxford}} & \multicolumn{3}{c}{\textbf{Severson}}\\
\cmidrule(lr){2-4}\cmidrule(l){5-7}
\textbf{Feature set} & \textbf{NASA src} & \textbf{CALCE src} & \textbf{ALL src} & \textbf{NASA src} & \textbf{CALCE src} & \textbf{ALL src}\\
\midrule
@@TAB_ABLATION@@
\bottomrule
\end{tabular}}
\end{table*}

\subsection{Failure-Definition Ablation: Measuring the Label Circularity}
Table~\ref{tab:faildef} relabels failure by a single endpoint at a time and repeats the transfer experiments. This quantifies the circularity concern of Section~II-A: with-SOH transfer \emph{tracks how much SOH defines the label}. Under the SOH-only label, transfer to Severson is @@FAILDEF_WITH_SOHLABEL@@; under the combined label, @@FAILDEF_WITH_COMBINED@@. Relabel by voltage sag alone and with-SOH transfer falls to @@FAILDEF_VOLT_WITH@@ --- a drop of about 0.17, with the with-minus-without gap narrowing from @@FAILDEF_ADV_COMBINED@@ to @@FAILDEF_ADV_VOLT@@. The residual (@@FAILDEF_VOLT_WITH@@ versus @@FAILDEF_VOLT_NO@@) is genuine but weak health-state information. Within one laboratory the distinction nearly vanishes --- on NASA the with-SOH model reaches @@FAILDEF_WITHIN_SOH@@ --- because sensor features co-move with SOH there; dataset shift destroys that co-movement, and with it the transfer illusion.

\begin{table*}[!t]
\caption{Failure-Definition Ablation at H=20 (XGBoost, ALL-LCO Source, Pooled AUC [CI]): Transfer With and Without SOH as a Feature, for Three Label Endpoints.}
\label{tab:faildef}
\centering
\resizebox{\textwidth}{!}{%
\begin{tabular}{lcccc}
\toprule
 & \multicolumn{2}{c}{\textbf{Oxford}} & \multicolumn{2}{c}{\textbf{Severson}}\\
\cmidrule(lr){2-3}\cmidrule(l){4-5}
\textbf{Label endpoint} & \textbf{with SOH} & \textbf{no SOH} & \textbf{with SOH} & \textbf{no SOH}\\
\midrule
@@TAB_FAILDEF@@
\bottomrule
\end{tabular}}
\end{table*}

\subsection{Same-Chemistry Cross-Laboratory Controls}
If the collapse of Section~III-C were a chemistry effect, transfer should survive when chemistry is held constant. Table~\ref{tab:samechem} shows it does not. NASA$\to$CALCE (both LCO) retains most of its transfer without SOH, so some laboratory pairs do share portable sensor signal; but Severson$\to$Oxford (both LFP) reproduces the full collapse pattern of the cross-chemistry shifts, including high with-SOH transfer and near-chance sensor-only transfer. These controls tighten the central attribution: the failure to transfer is a property of \emph{dataset shift} --- laboratories, protocols, and logging conventions --- that chemistry change adds to but does not uniquely cause. SOH is therefore a \emph{dataset- and chemistry-dependent} health-state signal, not a chemistry-specific one, and chemistry is not isolated as a causal variable.

\begin{table*}[!t]
\caption{Same-Chemistry Cross-Laboratory Controls at H=20 (XGBoost, Pooled AUC [CI]). Full Feature Sets Exercise the Zero-Imputation Confound; Common Sets Restrict Both Datasets to the Shared Four-Feature Subset.}
\label{tab:samechem}
\centering
\resizebox{\textwidth}{!}{%
\begin{tabular}{lccc}
\toprule
\textbf{Pair} & \textbf{with SOH (full)} & \textbf{no SOH (full)} & \textbf{no SOH (common)}\\
\midrule
@@TAB_SAMECHEM@@
\bottomrule
\end{tabular}}
\end{table*}

\subsection{Discrete-Time Hazard Model Versus Fixed-Horizon Classifiers}
Table~\ref{tab:hazard} compares the survival-product model of Section~II-G against the fixed-horizon XGBoost classifier. Within datasets the two match to within a few AUC points at every horizon (Table~\ref{tab:within}); no formal equivalence test is run. Under transfer the survival product inherits the same SOH dependence --- model formulation does not exempt inputs from dataset shift --- but $R(t,H)$ is monotone in $H$ by construction, whereas the four independent classifiers violate $p_{10}\le p_{20}\le p_{30}\le p_{50}$ on @@MONO_WITHIN_RANGE@@ of within-dataset rows (raw scores; Table~\ref{tab:mono}), rising to @@MONO_SEV_MAX@@ on transferred Severson rows. Monotone cumulative risk is what makes ``risk within 50 cycles'' coherent to quote: the hazard model buys structural coherence and a single fit at no within-dataset cost, and pays the same SOH shortcut as every other model.

\begin{table*}[!t]
\caption{Discrete-Time Hazard Model at H=20: Transfer Performance (Pooled AUC [CI]). The Fixed-Horizon XGBoost Comparator With SOH Is in Table~\ref{tab:crosswith}.}
\label{tab:hazard}
\centering
\resizebox{\textwidth}{!}{%
\begin{tabular}{llcc}
\toprule
\textbf{Training} & \textbf{Model} & \textbf{Oxford} & \textbf{Severson}\\
\midrule
@@TAB_HAZARD@@
\bottomrule
\end{tabular}}
\end{table*}

\begin{table}[!t]
\caption{Monotonicity Violations: Share of Scored Rows With $p_{10}\le p_{20}\le p_{30}\le p_{50}$ Failing, on Raw Scores (Platt in Parentheses). The Survival Product Is Monotone by Construction (Rate 0).}
\label{tab:mono}
\centering
\resizebox{\columnwidth}{!}{%
\begin{tabular}{lllcc}
\toprule
\textbf{Setting} & \textbf{Target} & \textbf{Model} & \textbf{Raw} & \textbf{Platt}\\
\midrule
@@TAB_MONO@@
\bottomrule
\end{tabular}}
\end{table}

\subsection{Failure of Calibration Transfer, and Its Partial Repair}
A second, independent failure mode appears in the cross-dataset tables when calibrated outputs are required. Isotonic regression is fitted as a step function on the source distribution; when target scores shift, many target scores fall into one bin and receive identical calibrated probabilities, and pooled ranking degrades accordingly --- in Table~\ref{tab:crosswith}, isotonic calibration costs the tree ensembles up to @@ISO_LOSS_OXF@@ pooled AUC on Oxford and @@ISO_LOSS_SEV@@ on Severson even with SOH present. Platt's smooth sigmoid degrades more gracefully, which is precisely the sense in which ``Platt better preserves discrimination at comparable calibration error.'' Fig.~\ref{fig:cal}(b) shows the deeper problem: on the shifted target, \emph{no} source-fitted calibrator lands near the diagonal, however it is fitted.

The recalibration arms of the companion study \cite{shikdar2026robustness} ask whether a small labeled target sample repairs this. Arm~A refits a calibrator on $k$ labeled target cells while leaving the classifier's scores untouched; Arm~B continues training the classifier on them. On Severson without SOH, zero-shot ECE is @@ARM_A_ECE0@@; at $k{=}5$ labeled cells, isotonic and Platt reduce it to @@ARM_A_ECE5_ISO@@ and @@ARM_A_ECE5_PLATT@@, while temperature scaling remains poorly calibrated (@@ARM_A_ECE5_TEMP@@). Calibration error is repaired, but discrimination is not: mean pooled AUC at $k{=}5$ stays near @@ARM_A_AUC5@@, because the calibrators create ties or reverse rankings on the shifted target. Arm~B is source-dependent: with NASA-trained XGBoost and $k{=}5$, Severson AUC rises from @@ARMB_NASA_XGB_ZERO@@ to @@ARMB_NASA_XGB_AFTER@@, overshooting the within-LCO reference ceiling (ratio @@ARMB_NASA_XGB_REC@@, an artifact of the mismatched ceiling rather than genuine super-ceiling recovery), while CALCE-trained updates leave the source operating point essentially untouched (LCO-holdout retention @@ARMB_CALCE_RET0@@ to @@ARMB_CALCE_RET1@@ for CALCE-trained XGBoost) and ALL-LCO updates damage it substantially (see Table~\ref{tab:rec-r2}). Tables~\ref{tab:rec-r1} and~\ref{tab:rec-r2} summarize both arms.

\begin{table*}[!t]
\caption{Arm A Recalibration on Severson Without SOH ($H{=}20$): Values Average the Three LCO Sources and Available Tree Models; ECE Uses 10 Equal-Width Probability Bins. From the Released Recalibration Study Artifacts.}
\label{tab:rec-r1}
\centering
\resizebox{\textwidth}{!}{%
\begin{tabular}{rccccccc}
\toprule
$k$ & Zero ECE & Iso. ECE & Platt ECE & Temp. ECE & Iso. AUC & Platt AUC & Temp. AUC\\
\midrule
@@TAB_REC_A@@
\bottomrule
\end{tabular}}
\end{table*}

\begin{table*}[!t]
\caption{Arm B Recovery at $k=5$ on Severson Without SOH, $H=20$. Recovery Ratio Uses the Within-LCO Ceiling; Retention Is the LCO-Holdout AUC Before and After the Update. \textsuperscript{*}Ratio above 1.0 is a ceiling-mismatch artifact, not super-ceiling recovery (see text).}
\label{tab:rec-r2}
\centering
\resizebox{\textwidth}{!}{%
\begin{tabular}{llccccc}
\toprule
\textbf{Training} & \textbf{Model} & \textbf{Zero AUC} & \textbf{Arm B AUC} & \textbf{Recovery} & \textbf{DeLong p} & \textbf{LCO retention}\\
\midrule
@@TAB_REC_B@@
\bottomrule
\end{tabular}}
\end{table*}

\subsection{Statistical Evidence for the SOH Ablation}
Table~\ref{tab:sohtests} states the central comparison with the cell as the experimental unit: the paired cell-level bootstrap interval on $\Delta$AUC $=$ AUC(with SOH) $-$ AUC(without SOH) for the ALL-LCO source, with cycle-level DeLong $p$-values as secondary evidence. On Severson, $\Delta$AUC reaches @@SOH_DELTA_SEV_XGB@@ (XGBoost, $H{=}20$) with the interval far from zero; on Oxford the point estimates agree but paired intervals are undefined, so Oxford contributes directional consistency only. DeLong values underflow throughout and are reported as $p{<}10^{-16}$ without interpretation (Section~\ref{sec:eval}). The claim rests on large bootstrapped effects on the 141-cell Severson target at every horizon, and on the failure-definition and same-chemistry controls of Sections~III-F and~III-G.

\begin{table*}[!t]
\caption{SOH Ablation Evidence (ALL-LCO Source): Pooled AUC With and Without SOH, Paired Cell-Level Bootstrap 95\% CI on the Difference (Primary), and Cycle-Level DeLong $p$-Value (Secondary; Inflated by Within-Cell Correlation).}
\label{tab:sohtests}
\centering
\resizebox{\textwidth}{!}{%
\begin{tabular}{llccccc}
\toprule
\textbf{Target} & \textbf{Model} & \textbf{H} & \textbf{AUC w/ SOH} & \textbf{AUC w/o} & \textbf{Delta AUC [boot 95\% CI]} & \textbf{DeLong p}\\
\midrule
@@TAB_SOH_TESTS@@
\bottomrule
\end{tabular}}
\end{table*}

\subsection{Operational Cost Analysis}
AUC treats all errors as equal; a BMS does not. We select the threshold on source out-of-fold scores with a 10\% false-alarm budget and apply it unchanged to the target (Table~\ref{tab:oper}). On Severson with SOH, the transferred XGBoost threshold misses @@OP_FNR_SEV_WITH@@ of failures at @@OP_FPR_SEV_WITH@@ false alarms. The same protocol fails for the other ensembles with SOH: LightGBM alerts on @@OP_FPR_SEV_WITH_LGBM@@ of negatives and Random Forest misses @@OP_FNR_SEV_WITH_RF@@ of failures --- threshold transfer fails for two of three ensembles even where ranking transfers. Without SOH, the XGBoost miss rate more than doubles to @@OP_FNR_SEV_NO@@ at @@OP_FPR_SEV_NO@@ false alarms; calibrated score types do not differ, because under shift the ranking is the asset and the ranking collapsed. On Oxford, nearly every row is positive, so no discriminating operating point exists there at all. The decision-curve view (Fig.~\ref{fig:netbenefit}) agrees: with SOH the model buys net benefit across thresholds; without SOH the curve hugs zero.

\begin{table*}[!t]
\caption{Operational Evaluation at H=20 (XGBoost, ALL-LCO Source): Decision Threshold Selected on Source Out-of-Fold Scores With a 10\% False-Alarm Budget, Then Applied Unchanged to the Target. Reported Are the Target Miss Rate (FNR) and Target False-Alarm Rate (FPR).}
\label{tab:oper}
\centering
\resizebox{\textwidth}{!}{%
\begin{tabular}{llcccccccc}
\toprule
 & & \multicolumn{2}{c}{\textbf{raw}} & \multicolumn{2}{c}{\textbf{Platt}} & \multicolumn{2}{c}{\textbf{isotonic}} & \multicolumn{2}{c}{\textbf{temperature}}\\
\cmidrule(lr){3-4}\cmidrule(l){5-6}\cmidrule(l){7-8}\cmidrule(l){9-10}
\textbf{Target} & \textbf{Features} & \textbf{FNR} & \textbf{FPR} & \textbf{FNR} & \textbf{FPR} & \textbf{FNR} & \textbf{FPR} & \textbf{FNR} & \textbf{FPR}\\
\midrule
@@TAB_OPER@@
\bottomrule
\end{tabular}}
\end{table*}

\begin{figure*}[!t]
\centering
\includegraphics[width=0.5\textwidth]{fig_netbenefit.png}
\caption{Decision-curve analysis (net benefit vs.\ threshold) for ALL-LCO $\to$ Severson without SOH at $H{=}20$: no score type produces a usable operating point once SOH is removed.}
\label{fig:netbenefit}
\end{figure*}

\subsection{SHAP: The Shortcut Is Visible in the Model}
TreeSHAP \cite{lundberg2020from} gives a direct view of the mechanism. Figs.~\ref{fig:shap_xgb}--\ref{fig:shap_rf} show the attributions for the three tree models in NASA-to-Oxford transfer at H=20 with SOH kept: SOH dominates by a wide margin, cycle index a distant second, and the voltage, current, temperature, and duration features contribute little. Figs.~\ref{fig:shap_noxgb}--\ref{fig:shap_norf} repeat the same comparisons with SOH removed, and every remaining feature collapses to near-zero SHAP spread --- matching, in the model's internal accounting, the quantitative collapse of Sections~III-C--III-E and the sensor-inversion of Section~III-D.

\begin{figure*}[!t]
\centering
\includegraphics[width=0.85\textwidth]{Fig06a_XGBoost_SHAP.png}
\caption{SHAP feature importance for XGBoost in NASA-to-Oxford transfer at H=20, with SOH. SOH dominates by a wide margin.}
\label{fig:shap_xgb}
\end{figure*}

\begin{figure*}[!t]
\centering
\includegraphics[width=0.85\textwidth]{Fig06b_LightGBM_SHAP.png}
\caption{SHAP feature importance for LightGBM in NASA-to-Oxford transfer at H=20, with SOH. SOH dominates by a wide margin.}
\label{fig:shap_lgbm}
\end{figure*}

\begin{figure*}[!t]
\centering
\includegraphics[width=0.85\textwidth]{Fig06c_RandomForest_SHAP.png}
\caption{SHAP feature importance for Random Forest in NASA-to-Oxford transfer at H=20, with SOH. SOH dominates by a wide margin.}
\label{fig:shap_rf}
\end{figure*}

\begin{figure*}[!t]
\centering
\includegraphics[width=0.85\textwidth]{Fig06d_XGBoost_SHAP_noSOH.png}
\caption{SHAP feature importance for XGBoost in NASA-to-Oxford transfer at H=20, without SOH. All features collapse to near-zero spread, so no transferable signal remains.}
\label{fig:shap_noxgb}
\end{figure*}

\begin{figure*}[!t]
\centering
\includegraphics[width=0.85\textwidth]{Fig06e_LightGBM_SHAP_noSOH.png}
\caption{SHAP feature importance for LightGBM in NASA-to-Oxford transfer at H=20, without SOH. All features collapse to near-zero spread.}
\label{fig:shap_nolgbm}
\end{figure*}

\begin{figure*}[!t]
\centering
\includegraphics[width=0.85\textwidth]{Fig06f_RandomForest_SHAP_noSOH.png}
\caption{SHAP feature importance for Random Forest in NASA-to-Oxford transfer at H=20, without SOH. All features collapse to near-zero spread.}
\label{fig:shap_norf}
\end{figure*}

\subsection{Embedded Deployment Validation}
Three independent validation stages confirm that the C tree engine reproduces the Python training libraries through the output transforms of \eqref{eq:rf} and \eqref{eq:sigmoid}: a Python walker over the extracted trees, the same walker cross-checked on the full 1~028-row set, and the C engine run on-device on an ESP32-S3 against the library's \texttt{predict\_proba()}. Fig.~\ref{fig:pyc} plots every one of the 1~028 on-device predictions against the Python values for all three models; the points fall on the identity line and the residual histograms are centered at zero. The measured on-device maximum absolute error is $2.4\times10^{-7}$ for XGBoost, $1.1\times10^{-6}$ for LightGBM, and $3.0\times10^{-8}$ for Random Forest, all far inside the $10^{-5}$ tolerance.

Table~\ref{tab:deploy} lists sizes, validation error, and inference time; Fig.~\ref{fig:deploy} plots latency and footprint. The ensembles occupy about 372 kB. At 240 MHz: from flash, 0.76, 1.55, and 2.36 ms per model; from PSRAM, 0.50, 0.93, and 1.29 ms; single-model from SRAM, 0.32, 0.43, and 0.50 ms.
The full 372 kB binary does not fit internal SRAM (free heap roughly 320 kB), but each model does, and the firmware runs one model at a time. Timings land above the 150--250 \textmu s projection because the working set exceeds the ESP32-S3's data cache; even so, the all-three PSRAM path clears the 100 ms per-cycle budget by roughly 37\texttimes. The deployment result is deliberately independent of the validity question: whatever the models are reading, on this dataset they can read it in under a millisecond on a \$12 microcontroller. The on-device outputs are raw ensemble probabilities through \eqref{eq:rf}--\eqref{eq:sigmoid} with no calibration layer: the Platt and isotonic maps live off-device, and Section~III-I shows source-fitted maps do not survive dataset shift, so a deployed system must revalidate or refit its calibrator on target-labeled cells rather than trust the laboratory calibration.

\begin{table*}[!t]
\caption{Embedded Deployment: Model Size, Accuracy, On-Device Validation Error, and Measured Inference Time. Errors Are Maximum $|$C engine $-$ \texttt{predict\_proba()}$|$ Over All 1 028 Rows. Timings Measured On-Device at 240 MHz With the Binary in Flash (PROGMEM), PSRAM, or a Single Model in Internal SRAM.}
\label{tab:deploy}
\centering
\resizebox{\textwidth}{!}{%
\begin{tabular}{lccccc}
\toprule
\textbf{Model} & \textbf{Nodes} & \textbf{Max Err} & \textbf{Flash} & \textbf{PSRAM} & \textbf{SRAM (1 model)}\\
\midrule
XGBoost & 3 272 & $2.4\times10^{-7}$ & 0.76 ms & 0.50 ms & 0.32 ms\\
LightGBM & 7 000 & $1.1\times10^{-6}$ & 1.55 ms & 0.93 ms & 0.43 ms\\
Random Forest & 16 064 & $3.0\times10^{-8}$ & 2.36 ms & 1.29 ms & 0.50 ms\\
All three & 26 336 & --- & 4.7 ms & 2.7 ms & ---\\
\bottomrule
\end{tabular}}
\end{table*}

\begin{figure*}[!t]
\centering
\includegraphics[width=0.9\textwidth]{fig_deploy.png}
\caption{Deployment trade-offs: (a) measured single-row inference time per model at 240 MHz with the binary in memory-mapped flash, PSRAM, or a single model in internal SRAM (dashed line: all three models from PSRAM); (b) per-model and combined binary footprint as a share of the ESP32-S3's roughly 320 kB of free internal SRAM --- each model fits, the combined 372 kB binary does not.}
\label{fig:deploy}
\end{figure*}

\begin{figure*}[!t]
\centering
\includegraphics[width=0.9\textwidth]{fig_py_c_agreement.png}
\caption{On-device C engine predictions versus Python \texttt{predict\_proba()} over all 1 028 rows and three models (agreement scatter on the identity line, inset residual histogram of C minus Python in units of $10^{-6}$). Maximum absolute errors are $2.4\times10^{-7}$ (XGBoost), $1.1\times10^{-6}$ (LightGBM), and $3.0\times10^{-8}$ (Random Forest).}
\label{fig:pyc}
\end{figure*}

\subsection{Replication on New Chemistries: NCA, NMC, and LFP Transfer Targets}
The transfer claim is re-tested on @@BA_CELLS_TOTAL@@ cells of NMC, NCA and LFP chemistry from two further laboratories (Battery Archive groups 5--8 of Section~II-B), with the LCO training pool and protocol untouched --- the new cells never enter model selection --- and the experiment run on the same protocol version as every result above. Fig.~\ref{fig:baquality} shows the post-cleaning capacity traces with the pre-registered audit decisions. Table~\ref{tab:crosschemnew} reports the NASA+CALCE $\rightarrow$ new-target transfer at $H{=}20$, with and without SOH, together with the paired cell-level bootstrap for the SOH ablation.

\begin{figure*}[!t]
\centering
\includegraphics[width=0.95\textwidth]{fig_ba_quality.png}
\caption{Data audit of the Battery Archive groups after the pre-registered cleaning rule (formation-cycle drop, Hampel despiking, level-shift truncation, voltage validity, minimum 20 usable cycles): (a) admitted-cell capacity traces with the 80\% end-of-life criterion; (b) retained versus excluded cycles per group; (c) recorded exclusion reasons. Of @@BA_CELLS_AUDIT@@ downloaded cells, @@BA_CELLS_TOTAL@@ are admitted and @@BA_CELLS_EXCLUDED@@ are excluded for exactly one recorded reason.}
\label{fig:baquality}
\end{figure*}

With SOH available, transfer from the LCO sources spans pooled AUC @@XFER_BA_WITH@@ across the new targets --- the same upper-bound regime as Table~\ref{tab:crosswith}. Removing SOH drops transfer to @@XFER_BA_NO@@, and the paired per-cell bootstrap puts the with-minus-without difference at @@XFER_BA_DELTA@@ AUC across the three ensembles --- on the SNL groups the LCO-trained ensemble contributes nothing portable on NCA, NMC, or LFP chemistry from an unseen laboratory, exactly as it contributed nothing on the earlier LFP targets. The HNEI NMC group is the exception: full no-SOH transfer (cycle index included) stays at 0.99 pooled there. All 15 HNEI cells cycle under one condition and one fixed charge protocol, so cycle-index and voltage cues survive the laboratory change; this is the dataset-shift mechanism of Section~III-G operating in reverse --- narrow, stereotyped protocols transfer, diverse ones do not --- not chemistry-specific knowledge. The common-set column repeats the check on the unit-safe feature subset (Section~II-C), controlling for the structurally missing current and temperature columns of the export: the SNL collapse is not an artifact of the zero-imputed columns --- though only for the available telemetry, per the Limitations. On the new groups the voltage-sag endpoint never fires before the SOH endpoint (@@BA_VOLT_FIRST@@ of @@BA_CELLS_TOTAL@@ admitted cells), so the composite label is, effectively, the SOH criterion alone; the failure-definition ablation of Section~III-F therefore cannot be re-run on these targets, and all with-SOH numbers here are to be read as label-proximity upper bounds, as defined in Section~II-A.

\begin{table*}[!t]
\caption{Transfer to New Chemistries at H=20 (NASA+CALCE Source). Pooled AUC With Cell-Level Bootstrap 95\% CI; $\Delta$ Is the Paired Per-Cell Bootstrap Difference (With $-$ Without SOH). New Groups Are Transfer Targets Only.}
\label{tab:crosschemnew}
\centering
\resizebox{\textwidth}{!}{%
\begin{tabular}{llccc}
\toprule
\textbf{Target} & \textbf{Model} & \textbf{With SOH} & \textbf{Without SOH} & \textbf{Delta (paired)}\\
\midrule
@@TAB_CROSSCHEM_NEW@@
\bottomrule
\end{tabular}}
\end{table*}

\subsection{Operating-Condition Shift at Fixed Chemistry and Laboratory}
The same-chemistry controls of Section~III-G varied laboratory while holding chemistry fixed; here the complement is measured --- hold chemistry \emph{and} laboratory fixed and vary only the operating condition. For each SNL group (NCA, NMC, LFP), the test cell is chosen from an operating condition (temperature band, C-rate, or both) that contributes no training cell in that fold; metadata parsed from cell identifiers provides the condition labels and is never used as a feature. Table~\ref{tab:condshift} compares this leave-condition-out protocol with the ordinary mixed-cell folds on the same group.

\begin{table}[!t]
\caption{Leave-Condition-Out Versus Mixed-Cell Folds at H=20 (XGBoost). Conditions Are Temperature Band, C-Rate, or Both; Metadata Never Enters the Features. Ref.\ Is the Ordinary Mixed-Cell AUC for the Same Feature Set (--- Where Unstable or Not Run).}
\label{tab:condshift}
\centering
\resizebox{\columnwidth}{!}{%
\begin{tabular}{llcccc}
\toprule
\textbf{Group} & \textbf{Held-out axis} & \textbf{With SOH} & \textbf{Ref.} & \textbf{No SOH} & \textbf{Ref.}\\
\midrule
@@TAB_CONDSHIFT@@
\bottomrule
\end{tabular}}
\end{table}

Holding out an entire operating condition at fixed chemistry and laboratory costs little with SOH and more without it: on SNL-NCA the with-SOH pooled AUC moves only from @@COND_NCA_REF@@ in ordinary mixed-cell folds to @@COND_NCA_TEMPRATE@@ under the joint holdout, and across all interpretable group--axis combinations the largest with-SOH drop is @@COND_DROP_WITH_MAX@@ AUC points --- whereas the sensor-only model degrades before the with-SOH one does, falling to @@COND_NCA_TEMPRATE_NO@@ on NCA (both axes) and @@COND_NMC_TEMPRATE_NO@@ on NMC. Two caveats bound the claim's strength. First, the mixed-fold raw reference on the 6-cell SNL-NMC group is itself unstable --- with six cells the split composition dominates, and raw AUC there (@@COND_NMC_MIXED@@ with SOH) disagrees with the @@BA_WITHIN_NMC_MIN@@ that the same model reaches under the paper's cell-disjoint outer folds --- so its reference column is marked --- and no drop is claimed for that group. Second, and most importantly, the distance-rule baseline is again undisturbed: its pooled AUC under condition holdout spans @@COND_DIST_MIN@@--@@COND_DIST_MAX@@, essentially its mixed-fold value, so whatever the rule keys on is invariant to operating condition. Two conclusions follow. The residual within-laboratory signal that survives dataset shift is only mildly condition-dependent --- chemistry and laboratory shift, not protocol condition, carry the larger penalty --- and no competing shortcut of comparable size hides in temperature or C-rate: with-SOH scores remain well above their no-SOH counterparts under condition holdout, so the SOH coordinate, not an operating-condition coordinate, is the portable quantity --- portable for the definitional reason quantified in Section~III-F. The HNEI group is omitted because its cells cycle under a single condition (25~$^\circ$C, mixed fixed discharge C-rates under one charge protocol), leaving no condition axis to hold out. SNL LFP was run with the with-SOH and common no-SOH feature sets only, so its full-feature no-SOH column reads --- by design, not by omission.

\subsection{Key Findings}
\begin{itemize}
\item Within-distribution models are reliable (fold-mean AUC @@WITHIN_MIN@@, Platt-calibrated) and run on-device in under a millisecond per SRAM-resident model.
\item Cross-dataset transfer collapses without SOH (Severson up to @@XFER_WITH_SEV@@ with SOH to @@XFER_NO_SEV_MIN@@--@@XFER_NO_SEV_MAX@@ without); an unfitted SOH distance rule matches the ensembles.
\item With-SOH numbers are diagnostic-proximity upper bounds: prognosis-only evaluation gives @@PROG_SEV@@ on Severson and @@PROG_BA@@ on the new targets.
\item Dataset and laboratory shift govern the collapse at least as much as chemistry; narrow protocols transfer, diverse ones do not, with no temperature or C-rate shortcut in the available telemetry.
\item Source-fitted calibration and transferred thresholds do not survive shift; target recalibration repairs ECE but not ranking.
\end{itemize}

\section{Discussion, Analysis, and Novelty}

\subsection{The SOH Shortcut, Precisely Stated}
The evidence assembled above supports a precise statement. When a failure-risk model trained on one battery dataset is evaluated on another, its apparent transfer performance is dominated by the SOH coordinate; much of that dominance is definitional, because SOH enters the failure label itself, and the failure-definition ablation of Section~III-F measures how much: relabeling failure by voltage sag alone removes @@FAILDEF_ADV_DROP@@ of the with-SOH transfer advantage on Severson. What remains is a weak, real health-state signal; everything else the sensors measure is dataset-specific, and the factorial ablation shows no non-SOH feature group carries usable cross-dataset signal. The replication of Section~III-N extends each of these statements to data the pipeline never had access to when it was designed. ``SOH is a chemistry-specific lookup key,'' the earlier framing, claimed more than the design could identify; the defensible claim is that apparent cross-dataset transferability is largely a health-state shortcut whose size this paper measures with label-, feature-, and dataset-level controls.

\subsection{Implications for Battery Management Practitioners}
Four practical points follow:
\begin{itemize}
\item \textbf{1)} Within-chemistry, within-laboratory models from standard charge--discharge features are reliable (fold-mean AUC above 0.89, Platt-calibrated) and run on a \$12 microcontroller in under a millisecond per SRAM-resident model. Deploy them where the training distribution holds.
\item \textbf{2)} Do not transfer models across datasets on the strength of an AUC table. If SOH is available at the target, an unfitted threshold rule is the correct first baseline; if SOH is unavailable, expect chance-level performance unless the target protocol is as narrow as the source.
\item \textbf{3)} Quote risks from a coherent formulation: four independent classifiers can rank ``risk within 10 cycles'' above ``risk within 50 cycles,'' while the survival product of Section~II-G cannot.
\item \textbf{4)} Under shift, prefer raw or Platt-scaled scores, expect calibration to need labeled target data, and choose thresholds on source data with an explicit cost model --- then verify.
\end{itemize}

\subsection{Limitations}
The honesty of this paper has limits of its own. Oxford contributes five cells --- one with a single positive row --- at horizon-invariant 100-cycle logging resolution, so it is a pooled-detection target re-verified on Severson. CALCE contributes seven cells, so its leave-one-cell-out uncertainty is wide. The Battery Archive groups sharpen the chemistry claim but not the condition claim: SNL NCA/NMC span few conditions with as few as three cells each, and mostly-censored SNL LFP is supporting evidence. The export lacks per-cycle current and temperature; affected features are zero-imputed, conclusions rest on the common feature set, and the sensor collapse cannot rule out richer-telemetry signals. The small three-seed GRU convicts no architecture beyond this instance, was not re-evaluated on the Battery Archive targets, and might differ from pretrained, transformer, or state-space models. The hazard model scores post-failure rows like every formulation. Cost analysis fixes 20:1 with source-selected thresholds. Transfer remains one-directional (LCO sources); isolating chemistry needs same-laboratory cross-chemistry campaigns.

\section{Conclusion}
This study widened a multi-horizon failure-risk framework and found the widening demanded a validity analysis. Within dataset the framework is sound: fold-mean AUC @@WITHIN_MIN@@ under cell-disjoint evaluation, a monotone discrete-time hazard formulation, and an about-372~kB deployment reproducing the training libraries to $1.1\times10^{-6}$ on a \$12 ESP32-S3.

Across datasets, the apparent transferability is mostly a State-of-Health shortcut. A distance-to-threshold rule on SOH alone reproduces the tuned ensembles' transfer performance; removing SOH collapses every model class toward or below chance; relabeling failure by voltage sag alone --- a criterion SOH does not define --- removes about a fifth of the with-SOH advantage; and same-chemistry cross-laboratory controls show the collapse is caused by dataset shift at least as much as by chemistry. The supported conclusion is deliberately narrower than ``SOH is chemistry-specific'': much of the observed cross-dataset discrimination is attributable to SOH, partly because SOH defines the failure criterion itself. Crucially, these statements replicate on @@BA_CELLS_TOTAL@@ NMC, NCA and LFP cells from two further laboratories: with-SOH transfer spans @@XFER_BA_WITH@@ and drops to @@XFER_BA_NO_SNL@@ without SOH outside the single-condition HNEI group (paired deltas @@XFER_BA_DELTA_SNL@@) --- a pre-registered replication, not a post-hoc fit.

Calibration behaves as the validity story predicts: Platt preserves discrimination best at comparable calibration error, isotonic fails under shift, labeled-target recalibration repairs calibration error but not discrimination, and no score type yields a usable operating point without SOH. Evidence is carried by cell-level bootstrap intervals, with cycle-level tests demoted to a consistency check.

Future work, in order of directness: an SOH-free failure definition as the primary cross-dataset target; same-laboratory, cross-chemistry aging campaigns; learned dataset-invariant representations; cell-level hierarchical calibration; and multi-seed, multi-architecture sequence-model studies at generalizable scale.

\clearpage
\section*{Acknowledgment}
This work received no external funding. The authors declare no conflicts of interest. The source code, data, trained models, firmware, and validation pipeline are publicly available at \url{https://github.com/touhidsiddiqueeraj-bit/Multi-Horizon-Hazard-Models-for-Battery-Failure-Prediction}.

\raggedbottom
\begin{thebibliography}{99}
\bibitem{shikdar2026learning} T. A. Shikdar and H. Laaksonen, ``Learning when not to use a battery: Multihorizon failure intelligence,'' \emph{Int. Trans. Electr. Energy Syst.}, vol. 36, Art. no. 6000810, 2026.
\bibitem{batteryarchive2023} Battery Archive, ``Open battery degradation testing data,'' 2023. [Online]. Available: \url{https://batteryarchive.org}. Consolidated cell cycling records from the Hawaii Natural Energy Institute and Sandia National Laboratories.
\bibitem{meng2019review} H. Meng and Y. Li, ``A review on prognostics and health management (PHM) methods of lithium-ion batteries,'' \emph{Renew. Sustain. Energy Rev.}, vol. 116, Art. no. 109405, 2019.
\bibitem{breiman2001random} L. Breiman, ``Random forests,'' \emph{Mach. Learn.}, vol. 45, no. 1, pp. 5--32, 2001.
\bibitem{chen2016xgboost} T. Chen and C. Guestrin, ``XGBoost: A scalable tree boosting system,'' in \emph{Proc. ACM SIGKDD}, 2016, pp. 785--794.
\bibitem{ke2017lightgbm} G. Ke, Q. Meng, T. Finley, et al., ``LightGBM: A highly efficient gradient boosting decision tree,'' in \emph{Proc. NeurIPS}, vol. 30, 2017, pp. 3146--3154.
\bibitem{saha2007battery} B. Saha and K. Goebel, ``Battery data set,'' NASA Ames Prognostics Data Repository, 2007.
\bibitem{calce2023battery} CALCE Battery Research Group, ``Battery aging datasets,'' Univ. Maryland, 2023.
\bibitem{birkl2017oxford} C. R. Birkl and D. A. Howey, ``Oxford battery degradation dataset 1,'' Univ. Oxford, 2017.
\bibitem{severson2019data} K. A. Severson, P. M. Attia, N. Jin, et al., ``Data-driven prediction of battery cycle life before capacity degradation,'' \emph{Nature Energy}, vol. 4, pp. 383--391, 2019.
\bibitem{sahoo2022transfer} S. Sahoo, K. S. Hariharan, S. Agarwal, et al., ``Transfer learning based generalized framework for state of health estimation of Li-ion cells,'' \emph{Sci. Rep.}, vol. 12, Art. no. 13173, 2022.
\bibitem{lu2023deep} J. Lu, R. Xiong, J. Tian, C. Wang, and F. Sun, ``Deep learning to estimate lithium-ion battery state of health without additional degradation experiments,'' \emph{Nat. Commun.}, vol. 14, Art. no. 2760, 2023.
\bibitem{platt1999probabilistic} J. C. Platt, ``Probabilistic outputs for support vector machines and comparisons to regularized likelihood methods,'' in \emph{Adv. Large Margin Classifiers}, MIT Press, 1999, pp. 61--74.
\bibitem{zadrozny2002transforming} B. Zadrozny and C. Elkan, ``Transforming classifier scores into accurate multiclass probability estimates,'' in \emph{Proc. ACM SIGKDD}, 2002, pp. 694--699.
\bibitem{niculescu2005predicting} A. Niculescu-Mizil and R. Caruana, ``Predicting good probabilities with supervised learning,'' in \emph{Proc. ICML}, 2005, pp. 625--632.
\bibitem{huang2020experimental} L. Huang, J. Zhao, B. Zhu, H. Chen, and S. K. L. M. v. Broucke, ``An experimental investigation of calibration techniques for imbalanced data,'' \emph{IEEE Access}, vol. 8, pp. 127343--127352, 2020.
\bibitem{lundberg2020from} S. M. Lundberg, G. G. Erion, H. Chen, et al., ``From local explanations to global understanding with explainable AI for trees,'' \emph{Nat. Mach. Intell.}, vol. 2, no. 1, pp. 56--67, 2020.
\bibitem{grzesik2024combining} P. Grzesik and D. Mrozek, ``Combining machine learning and edge computing: Opportunities, challenges, platforms, frameworks, and use cases,'' \emph{Electronics}, vol. 13, no. 3, Art. no. 640, 2024.
\bibitem{gupta2023online} C. Gupta and A. Ramdas, ``Online Platt scaling with calibeating,'' in \emph{Proc. ICML}, PMLR 202, 2023, pp. 12182--12204.
\bibitem{preger2020degradation} Y. Preger, H. C. Barkholtz, A. Fresquez, et al., ``Degradation of commercial lithium-ion cells as a function of chemistry and cycling conditions,'' \emph{J. Electrochem. Soc.}, vol. 167, no. 12, Art. no. 120532, 2020.
\bibitem{attia2020closedloop} P. M. Attia, A. Grover, N. Jin, et al., ``Closed-loop optimization of fast-charging protocols for batteries with machine learning,'' \emph{Nature}, vol. 578, no. 7795, pp. 397--402, 2020.
\bibitem{sulzer2021challenge} V. Sulzer, P. Mohtat, A. Aitio, et al., ``The challenge and opportunity of battery lifetime prediction from field data,'' \emph{Joule}, vol. 5, no. 8, pp. 1934--1955, 2021.
\bibitem{roman2021machine} D. Roman, S. Saxena, V. Robu, et al., ``Machine learning pipeline for battery state-of-health estimation,'' \emph{Nat. Mach. Intell.}, vol. 3, no. 4, pp. 447--456, 2021.
\bibitem{delong1988comparing} E. R. DeLong, D. M. DeLong, and D. L. Clarke-Pearson, ``Comparing the areas under two or more correlated receiver operating characteristic curves: A nonparametric approach,'' \emph{Biometrics}, vol. 44, no. 3, pp. 837--845, 1988.
\bibitem{shikdar2026robustness} T. A. Shikdar and H. Laaksonen, ``How far does the model hold? Robustness of multihorizon battery failure intelligence under realistic operating conditions,'' submitted for publication, 2026.
\end{thebibliography}

\EOD

\end{document}
